\begin{titlepage}
\thispagestyle{empty}
\vspace*{0.65cm}
\begin{center}
{\Huge\mdseries 大数据与管理决策基础\par}
\vspace{0.75cm}
{\Large 数据分析、统计推断与管理决策\par}
\vspace{0.5cm}
{\small\sffamily\color{BDMGray} Big Data and Management Decision Foundations\par}
\end{center}
\vspace{1.1cm}
\begin{center}
\includegraphics[width=\textwidth]{assets/cover_data_decisions.png}
\end{center}
\vspace{0.7cm}
\begin{center}
{\Large 孙振宇\quad 著\par}
\vspace{0.55cm}
{\large 东北大学数学与统计学院\par}
\vfill
{\small\color{BDMGray} 数智化企业运营与优化微专业\par}
\end{center}
\end{titlepage}

\chapter*{前言}
\addcontentsline{toc}{chapter}{前言}
\markboth{前言}{前言}
\begingroup
\setlength{\parskip}{0pt}

企业在生产、交易和服务过程中积累了大量数据。订单记录反映需求，设备日志记录
生产状态，财务资料衡量资源消耗与经营成果。这些资料为管理决策提供了观察基础。
从观察到决策，还需要明确记录的含义，评估统计结论的不确定性，并比较不同方案的
预期后果。本书围绕这一过程，介绍企业数据管理、统计分析与管理科学的基本理论
及其在生产运营中的应用。

数据管理、统计分析和组织治理各有研究对象，也相互依赖。数据模型规定对象的
身份、属性与关系，统计方法研究从样本获得总体知识的条件，治理制度规定决策权、
执行权与监督责任。把三者结合起来，才能使指标具有明确含义，使推断建立在适当的
假设上，使决策符合组织目标与资源条件。

全书分为三个部分。第一部分讨论决策与数据基础。第1章定义状态、行动、损失和
信息价值；第2至4章依次讨论企业数据模型、指标查询和数据质量。第二部分讨论
统计分析与效果评价。第5章研究分布、波动和排队规律，第6至8章分别研究随机实验、
预测和观测性因果推断。第三部分讨论优化与执行。第9章在成本、资源和风险约束下
选择方案，第10章通过事件日志研究业务流程，并讨论智能体参与执行的条件。

各章按照问题、定义、假设、推导和应用展开。信息价值由获得信息前后的最优期望
损失之差定义；实验估计的方差由处理分配机制决定；预测区间的覆盖性质由相应的
概率条件保证；优化结果由目标函数和可行域共同确定。公式之后的算例用于说明
计算步骤，条件变化的讨论用于说明结论如何随之改变。

企业应用取材于论文、官方技术文档和企业披露，包括 Palantir、Oracle、Databricks、
Atlan、Collibra、Monte Carlo 和 Celonis 等案例。来源引文对应公开事实，随后的
模型与数值算例用于分析其技术机制。除注明来源的数据外，算例数值均为教学设定。
制造案例使用需求、设备、质量与交付等十三类分析要素组织材料，其定义在第2章给出。

人工智能应用涉及查询、预测与执行。企业系统分析从计算资源、数据、业务
语义、模型服务、任务运行和实际行动出发，分别说明其技术条件与管理责任。
自然语言查询需要确定指标含义与数据范围，检索增强生成需要取得适用的资料，
智能体执行需要管理工具权限与任务状态。各章将这些要求落实为查询检验、
统计评价、资源核算和状态转换规则。

本书面向具有初等微积分、线性代数和概率统计基础的本科生及企业实践人员。
符号在首次使用时定义，附录汇集数学关系、计算复现约定和证据写作规范。
正文保留必要代码，完整运行入口集中于附录；参考文献与主题索引统一置于书末。

\begin{flushright}
孙振宇\qquad 2026年9月
\end{flushright}
\endgroup

\chapter*{全书结构与记号}
\addcontentsline{toc}{chapter}{全书结构与记号}
\markboth{全书结构与记号}{全书结构与记号}

\section*{全书的组织}
图~\ref{fig:book-map}~概括各部分的知识联系。问题定义确定所需信息，数据模型与
查询形成分析资料，统计方法提供分布和效果估计，优化方法据此比较方案。
执行结果又成为后续分析的观测资料。数据质量检查和目标修订贯穿这一过程。

\begin{figure}[htbp]
\centering
\begin{tikzpicture}[x=1cm,y=1cm]
\node[bdmnode,text width=3.1cm] (a) at (0,0) {第1章\\决策问题与损失};
\node[bdmnode,text width=3.1cm] (b) at (4.4,0) {第2--4章\\对象、指标与可信度};
\node[bdmnode,text width=3.1cm] (c) at (8.8,0) {第5章\\分布、波动与运营};
\node[bdmnode,text width=3.1cm] (d) at (8.8,-2) {第6--8章\\实验、预测与因果};
\node[bdmnode,text width=3.1cm] (e) at (4.4,-2) {第9章\\约束下的行动选择};
\node[bdmnode,text width=3.1cm] (f) at (0,-2) {第10章\\执行、记录与反馈};
\draw[bdmarrow] (a)--(b);\draw[bdmarrow] (b)--(c);
\draw[bdmarrow] (c)--(d);\draw[bdmarrow] (d)--(e);
\draw[bdmarrow] (e)--(f);\draw[bdmarrow,dashed] (f)--(a);
\end{tikzpicture}
\caption{决策问题、数据基础、统计分析与方案执行的关系。}
\label{fig:book-map}
\end{figure}

第6章的随机实验和第8章的观测性研究讨论同一类处理效应，区别在于处理分配机制。
第7章的预测模型则研究给定信息下的未知结果，为第9章提供需求或风险估计。
因此，预测与因果推断是两种不同的统计任务，各自的结论经过成本与资源分析后，
都可以用于方案选择。第10章把所选方案与实际流程联系起来，研究执行及其结果。

委托代理、置信序列、共形预测、双重稳健与尾部风险的推导分别依赖优化理论、
概率过程或统计推断的相应先修知识。

\section*{企业人工智能系统的分析范围}
企业需要先确定希望改善的业务结果，随后说明可观察的状态、可实施的行动和评价
窗口。计算资源和数据系统提供处理能力；业务对象与指标定义赋予记录以可供查询的
含义；模型服务完成预测、检索或生成；运行程序调度工具并保存状态；业务系统执行
获得授权的操作。访问控制、质量检查与责任分配作用于每一个环节。

这些职能可以由同一平台提供，也可以分布在数家供应商和企业自有系统中。因此，
技术图应同时标明数据的位置、调用接口、使用者身份、可写入的业务对象以及收费
单位。第2章研究数据与语义，第3、4章研究查询与质量，第5至8章研究运营度量及
效果评价，第9章研究资源约束，第10章研究任务执行与恢复。模型服务容量分析见
附录~\ref{app:ai-infra-lab}。在该算例中，计算出的吞吐与费用均为指定假设下的
估计值；投入使用前仍需测量真实负载和服务行为。

\section*{贯穿案例的使用约定}
本书反复使用一家假想封测工厂的订单履约问题。订单表示客户承诺，批次表示生产中
流动的对象，工序事件表示对象经过的活动，设备记录提供产能和运行状态。
一张订单可以拆成多个批次，一个批次可以多次经过同一设备，一次发运也可以包含
多张订单。各章按所研究的问题选择记录单位，并说明对象之间的对应关系。

随机实验另以结账页比较为例，预测以客户流失为例，优化以库存分配为例。
这些算例分别说明随机分配、标签定义和资源约束。相应方法应用于制造情境时，
须根据生产对象重新确定时间、处理和结果变量。第~\ref{sec:book-synthesis}~节
以订单加急问题综合运用各章方法。

\clearpage
\section*{数学与统计记号}
\begin{longtable}{P{0.22\linewidth}P{0.68\linewidth}}
\toprule
\rowcolor{BDMTableHead}
记号 & 约定与解释\\
\midrule
\endhead
\(X,x\) & 大写表示随机变量或特征向量，小写表示其实现值；状态已知时写为 \(X=x\)。\\
\(Y,Y(a)\) & \(Y\) 为结果或外生状态；\(Y(a)\) 强调结果会随行动改变。具体含义在各模型中声明。\\
\(a,\mathcal A(x)\) & 行动及状态 \(x\) 下的可行行动集合。\\
\(L,\ell,C\) & 业务损失、统计学习损失和另行计入的成本；同一成本不得重复计算。\\
\(\E,\Prob,\operatorname{Var},\operatorname{Cov}\) & 期望、概率、方差和协方差；条件对象写在竖线之后。\\
\(\theta,\hat\theta\) & 未知参数及其估计值；估计量是数据的函数，其抽样变异由标准误等量刻画。\\
\(n,N,i,t\) & 样本量、有限总体规模、对象索引和时间索引；局部例子可能重新定义并说明。\\
\(D,e(x),\tau(x)\) & 因果章节中分别表示处理、倾向得分和条件平均处理效应。\\
\(\mathbf1\{B\}\) & 事件 \(B\) 成立时为1，否则为0。\\
\(\argmin,\argmax\) & 使目标达到最小、最大的取值或取值集合；最优解可能不唯一。\\
\(\lambda\) & 局部定义的成本权重、正则化系数或到达率；不能跨模型混用量纲。\\
\bottomrule
\end{longtable}

等号用于定义或精确结论，近似号用于数值舍入、渐近近似或模型近似。除另有说明，
金额例子不含税费和折现，比例以0至1表示，百分点描述两个比例的绝对差。
例如从40\%提高到50\%，变化为10个百分点，相对增幅为25\%。对置信水平、显著性
水平、覆盖率和成功率，本书分别给出其概率所针对的对象，并在具体模型中说明其计算条件。
